%%
%% citemsp-doc.tex — Documentation for the citemsp package
%%
%% Copyright (c) 2025 Apostolos Tsampodimos
%% LPPL 1.3c — see LICENSE for details
%%
\documentclass[11pt, a4paper]{article}

\usepackage{lmodern}
\usepackage{microtype}
\usepackage[margin=2.5cm]{geometry}
\usepackage[dvipsnames, table]{xcolor}
\definecolor{linkblue}{rgb}{0.0, 0.2, 0.6}
\definecolor{linkred}{rgb}{0.7, 0.13, 0.13}
\usepackage[style=numeric, backend=biber]{biblatex}
\usepackage{citemsp}
\usepackage[colorlinks=true, citecolor=linkblue, linkcolor=linkred, urlcolor=linkblue]{hyperref}
\usepackage{tcolorbox}
\usepackage{authblk}
\usepackage{graphicx}
\usepackage{array}
\usepackage{booktabs}
\usepackage{orcidlink}
\usepackage{fontawesome5}
\usepackage{enumitem}

% ── Helpers ───────────────────────────────────────────────────────────────────
\newcommand{\githublink}[1]{\href{https://github.com/#1}{\faGithub}}
\newcommand{\pkg}[1]{\texttt{#1}}
\newcommand{\cmd}[1]{\texttt{\textbackslash #1}}

% ── Example bibliography ─────────────────────────────────────────────────────
\begin{filecontents}{refs-citemsp.bib}
@article{einstein1915,
  author  = {Einstein, Albert},
  title   = {Die Feldgleichungen der Gravitation},
  journal = {Sitzungsberichte der Preussischen Akademie der Wissenschaften},
  year    = {1915},
  pages   = {844--847}
}
@book{wald1984,
  author    = {Wald, Robert M.},
  title     = {General Relativity},
  publisher = {University of Chicago Press},
  year      = {1984}
}
@article{hawking1974,
  author  = {Hawking, Stephen W.},
  title   = {Black hole explosions?},
  journal = {Nature},
  volume  = {248},
  pages   = {30--31},
  year    = {1974}
}
@book{mtw1973,
  author    = {Misner, Charles W. and Thorne, Kip S. and Wheeler, John Archibald},
  title     = {Gravitation},
  publisher = {W. H. Freeman},
  year      = {1973}
}
\end{filecontents}
\addbibresource{refs-citemsp.bib}

% ── Title ─────────────────────────────────────────────────────────────────────
\title{\textbf{The \pkg{citemsp} Package}\\[4pt]
  \large Per-Key Citation Locators for Biblatex\\[4pt]
  \normalsize v2.1 — 2025/04/21}

\author{Apostolos Tsampodimos\,\,\githublink{atsampodimos/citemsp}\,\,\orcidlink{0009-0001-0230-5647}}
\affil{Center for Physical Sciences and Technology\\
Saul\.etekio~3, 10257 Vilnius, Lithuania\\
\texttt{apostolos.tsampodimos@ftmc.lt}}

\date{}

\begin{document}

\maketitle

\begin{abstract}
\noindent
The \pkg{citemsp} package extends \pkg{biblatex} with a single user-facing command, \cmd{citemsp}, that attaches section~(\S), paragraph~(\P), and prefix-based locators directly to numeric citation labels. An extensible prefix registry provides seven built-in types --- equation~(eq.), box~($\square$), definition~($\triangleq$), figure~($\boxdot$), footnote~($\dagger$), page~(pp.), and table~($\boxplus$) --- and users can register custom prefixes with \cmd{citemspprefix}. Standard \pkg{biblatex} provides locator fields through its postnote mechanism, but these cannot display stacked locator pairs per key in a compact inline form. The \pkg{citemsp} package solves this by parsing a slash-delimited syntax, rendering each locator as a superscript/subscript pair on the corresponding citation number, producing output such as~\citemsp{wald1984/3.2/5, mtw1973/b21.1, hawking1974/1/e4.3}. The package is compatible with all \pkg{biblatex} numeric styles and requires only \pkg{biblatex}, \pkg{graphicx}, and \pkg{amssymb}.
\end{abstract}

\tableofcontents

% ══════════════════════════════════════════════════════════════════════════════
\section{Motivation}
% ══════════════════════════════════════════════════════════════════════════════

The \pkg{biblatex} package offers a locator interface through its postnote field. A user can write \verb|\autocite[§3.2]{wald1984}| to append a free-form locator after the citation number. For multi-key citations, \verb|\autocites| accepts per-key pre- and post-notes:
%
\begin{tcolorbox}[colback=gray!5, colframe=gray!40, boxrule=0.5pt, arc=2pt, top=6pt, bottom=6pt, left=10pt, right=10pt]
\begin{verbatim}
\autocites[§3.2]{wald1984}[¶2]{hawking1974}
\end{verbatim}
\end{tcolorbox}
%
\noindent This approach has two limitations:

\begin{enumerate}[leftmargin=2em]
  \item \textbf{No stacked locators.} There is no built-in way to show both a section \emph{and} a paragraph locator for the same key in a single compact glyph. One must choose one or resort to verbose free text like ``§3.2, ¶5''.
  \item \textbf{Verbose syntax.} Citing three sources with locators requires nested optional arguments, which is cumbersome in source files.
  \item \textbf{No typed locators.} There is no compact notation for referencing specific equations, boxes, figures, tables, definitions, footnotes, or pages within a cited work.
\end{enumerate}

\noindent The \pkg{citemsp} package addresses all three: it renders locators as a superscript/subscript pair attached to the citation number, supports section and paragraph locators by default plus seven additional types via an extensible prefix registry, and accepts all keys with their locators in a single comma-separated argument.


% ══════════════════════════════════════════════════════════════════════════════
\section{Installation}
% ══════════════════════════════════════════════════════════════════════════════

\subsection{Per-project}

Copy \texttt{citemsp.sty} into the same directory as your main \texttt{.tex} file.

\subsection{System-wide (local texmf)}

\begin{tcolorbox}[colback=gray!5, colframe=gray!40, boxrule=0.5pt, arc=2pt, top=6pt, bottom=6pt, left=10pt, right=10pt]
\begin{verbatim}
mkdir -p ~/texmf/tex/latex/citemsp
cp citemsp.sty ~/texmf/tex/latex/citemsp/
texhash ~/texmf
\end{verbatim}
\end{tcolorbox}


% ══════════════════════════════════════════════════════════════════════════════
\section{Usage}
% ══════════════════════════════════════════════════════════════════════════════

\subsection{Loading}

The package must be loaded \emph{after} \pkg{biblatex}:
%
\begin{tcolorbox}[colback=gray!5, colframe=gray!40, boxrule=0.5pt, arc=2pt, top=6pt, bottom=6pt, left=10pt, right=10pt]
\begin{verbatim}
\usepackage[style=numeric, backend=biber]{biblatex}
\usepackage{citemsp}
\end{verbatim}
\end{tcolorbox}
%
\noindent Loading \pkg{citemsp} before \pkg{biblatex} produces a package error.


\subsection{Syntax}

The \cmd{citemsp} command accepts a comma-separated list of entries.  Each entry has the form \texttt{key}, \texttt{key/loc1}, or \texttt{key/loc1/loc2}. By default, the first slot renders as \S\ (section) and the second as \P\ (paragraph). A single-letter prefix changes the locator type:

\begin{tcolorbox}[colback=gray!5, colframe=gray!40, boxrule=0.5pt, arc=2pt, top=6pt, bottom=6pt, left=10pt, right=10pt]
\texttt{\textbackslash citemsp\{key\}} \hfill $\rightarrow$ \citemsp{wald1984}\\[3pt]
\texttt{\textbackslash citemsp\{key/section\}} \hfill $\rightarrow$ \citemsp{wald1984/6.2}\\[3pt]
\texttt{\textbackslash citemsp\{key/section/paragraph\}} \hfill $\rightarrow$ \citemsp{wald1984/6.2/3}\\[3pt]
\texttt{\textbackslash citemsp\{key/section/e4.3\}} \hfill $\rightarrow$ \citemsp{wald1984/6.2/e4.3}\\[3pt]
\texttt{\textbackslash citemsp\{key/b21.1\}} \hfill $\rightarrow$ \citemsp{mtw1973/b21.1}\\[3pt]
\texttt{\textbackslash citemsp\{key/section/b8.1\}} \hfill $\rightarrow$ \citemsp{mtw1973/8.4/b8.1}
\end{tcolorbox}

\noindent The built-in prefixes are: \texttt{b}~for box~($\square$), \texttt{d}~for definition~($\triangleq$), \texttt{e}~for equation~(eq.), \texttt{f}~for figure~($\boxdot$), \texttt{n}~for footnote~($\dagger$), \texttt{p}~for page~(pp.), and \texttt{t}~for table~($\boxplus$). No prefix defaults to \S\ or \P\ depending on slot position. When both locators are present, the first appears as a superscript and the second as a subscript. Users can register custom prefixes with \cmd{citemspprefix}:
%
\begin{tcolorbox}[colback=gray!5, colframe=gray!40, boxrule=0.5pt, arc=2pt, top=6pt, bottom=6pt, left=10pt, right=10pt]
\begin{verbatim}
\citemspprefix{w}{w.}   % register "w" prefix rendering as "w."
\end{verbatim}
\end{tcolorbox}


\subsection{Configuration}

The locator glyph size is controlled by \cmd{citemspscale} (default~0.35). Redefine it after loading:
%
\begin{tcolorbox}[colback=gray!5, colframe=gray!40, boxrule=0.5pt, arc=2pt, top=6pt, bottom=6pt, left=10pt, right=10pt]
\begin{verbatim}
\renewcommand{\citemspscale}{0.4}
\end{verbatim}
\end{tcolorbox}


% ══════════════════════════════════════════════════════════════════════════════
\section{Comparison with Standard Biblatex}
% ══════════════════════════════════════════════════════════════════════════════

\begin{table}[ht]
\centering
\rowcolors{2}{blue!3}{white}
\begin{tabular}{>{\raggedright\arraybackslash}p{0.46\textwidth} >{\raggedright\arraybackslash}p{0.46\textwidth}}
\toprule
\rowcolor{blue!15}
\textbf{Standard biblatex} & \textbf{citemsp} \\
\midrule
\verb|\autocite[§3.2]{wald1984}| produces ``[1, §3.2]'' --- locator follows the number as free text &
\verb|\citemsp{wald1984/3.2}| produces \citemsp{wald1984/3.2} --- locator attached as compact glyph \\[6pt]

No built-in stacked section + paragraph &
\verb|\citemsp{wald1984/3.2/5}| produces \citemsp{wald1984/3.2/5} with both \\[6pt]

Multi-cite needs \verb|\autocites| with repeated optional arguments &
Single command: \verb|\citemsp{k1/s1/p1, k2/s2}| \\[6pt]

No typed locators &
Seven built-in prefixes: \texttt{b}~($\square$), \texttt{d}~($\triangleq$), \texttt{e}~(eq.), \texttt{f}~($\boxdot$), \texttt{n}~($\dagger$), \texttt{p}~(pp.), \texttt{t}~($\boxplus$); extensible via \cmd{citemspprefix} \\[6pt]

Postnote applies after last key only (in \verb|\cite|) &
Each key carries independent locators \\
\bottomrule
\end{tabular}
\caption{Feature comparison between standard \pkg{biblatex} locators and \pkg{citemsp}.}
\label{tab:comparison}
\end{table}


% ══════════════════════════════════════════════════════════════════════════════
\section{Examples}
% ══════════════════════════════════════════════════════════════════════════════

\subsection{Single source}

\begin{tcolorbox}[colback=blue!3, colframe=blue!30, boxrule=0.5pt, arc=2pt, top=6pt, bottom=6pt, left=10pt, right=10pt]
\begin{verbatim}
The Schwarzschild metric~\citemsp{wald1984/6.1/2} describes
spherically symmetric vacuum solutions.
\end{verbatim}
\end{tcolorbox}

\noindent\textbf{Output:} The Schwarzschild metric~\citemsp{wald1984/6.1/2} describes spherically symmetric vacuum solutions.


\subsection{Multiple sources with mixed locators}

\begin{tcolorbox}[colback=blue!3, colframe=blue!30, boxrule=0.5pt, arc=2pt, top=6pt, bottom=6pt, left=10pt, right=10pt]
\begin{verbatim}
General covariance~\citemsp{einstein1915, wald1984/4,
  hawking1974/1/2} is the foundation of general relativity.
\end{verbatim}
\end{tcolorbox}

\noindent\textbf{Output:} General covariance~\citemsp{einstein1915, wald1984/4, hawking1974/1/2} is the foundation of general relativity.


\subsection{Italic context}

The package detects italic font shape and adjusts kerning automatically:

\begin{tcolorbox}[colback=blue!3, colframe=blue!30, boxrule=0.5pt, arc=2pt, top=6pt, bottom=6pt, left=10pt, right=10pt]
\begin{verbatim}
\textit{Hawking radiation~\citemsp{hawking1974/1/2} implies
black holes are not truly black.}
\end{verbatim}
\end{tcolorbox}

\noindent\textbf{Output:} \textit{Hawking radiation~\citemsp{hawking1974/1/2} implies black holes are not truly black.}


\subsection{Equation locator}

Use the \texttt{e} prefix to reference a specific equation:

\begin{tcolorbox}[colback=blue!3, colframe=blue!30, boxrule=0.5pt, arc=2pt, top=6pt, bottom=6pt, left=10pt, right=10pt]
\begin{verbatim}
The field equations~\citemsp{wald1984/3.2/e4.3} can be derived
from a variational principle.
\end{verbatim}
\end{tcolorbox}

\noindent\textbf{Output:} The field equations~\citemsp{wald1984/3.2/e4.3} can be derived from a variational principle.


\subsection{Box locator}

Use the \texttt{b} prefix for box environments, as used extensively in Misner, Thorne \& Wheeler's \textit{Gravitation}~\cite{mtw1973}:

\begin{tcolorbox}[colback=blue!3, colframe=blue!30, boxrule=0.5pt, arc=2pt, top=6pt, bottom=6pt, left=10pt, right=10pt]
\begin{verbatim}
The stress-energy tensor~\citemsp{mtw1973/b21.1} is discussed
in a self-contained box.
\end{verbatim}
\end{tcolorbox}

\noindent\textbf{Output:} The stress-energy tensor~\citemsp{mtw1973/b21.1} is discussed in a self-contained box.


\subsection{Mixed locator types}

All locator types can be freely mixed in a single call:

\begin{tcolorbox}[colback=blue!3, colframe=blue!30, boxrule=0.5pt, arc=2pt, top=6pt, bottom=6pt, left=10pt, right=10pt]
\begin{verbatim}
See~\citemsp{wald1984/6.1/2, mtw1973/b21.1, hawking1974/1/e4.3}
for further details.
\end{verbatim}
\end{tcolorbox}

\noindent\textbf{Output:} See~\citemsp{wald1984/6.1/2, mtw1973/b21.1, hawking1974/1/e4.3} for further details.


\subsection{Additional locator types}

The prefix registry provides seven built-in types beyond the default \S/\P:

\begin{tcolorbox}[colback=blue!3, colframe=blue!30, boxrule=0.5pt, arc=2pt, top=6pt, bottom=6pt, left=10pt, right=10pt]
\begin{verbatim}
\citemsp{wald1984/3.2/d5}          % definition
\citemsp{wald1984/6.1/f3}          % figure
\citemsp{hawking1974/1/n7}         % footnote
\citemsp{wald1984/p42}             % page
\citemsp{mtw1973/8.4/t2}          % table
\end{verbatim}
\end{tcolorbox}

\noindent\textbf{Output:}
\citemsp{wald1984/3.2/d5},\enspace
\citemsp{wald1984/6.1/f3},\enspace
\citemsp{hawking1974/1/n7},\enspace
\citemsp{wald1984/p42},\enspace
\citemsp{mtw1973/8.4/t2}.


% ══════════════════════════════════════════════════════════════════════════════
\section{Implementation}
% ══════════════════════════════════════════════════════════════════════════════

The package (\textasciitilde80 lines) consists of four components:

\begin{enumerate}[leftmargin=2em]
  \item \textbf{Locator scaling.} \cmd{citemspscale} (default~0.35) and \cmd{citesize} control the size of locator glyphs via \cmd{scalebox} from \pkg{graphicx}.

  \item \textbf{Prefix registry.} An extensible dispatch mechanism based on \verb|\ifcsname|/\verb|\csname| (e-TeX). Seven built-in prefixes are registered: \texttt{b}~($\square$), \texttt{d}~($\triangleq$), \texttt{e}~(eq.), \texttt{f}~($\boxdot$), \texttt{n}~($\dagger$), \texttt{p}~(pp.), \texttt{t}~($\boxplus$). Users can add custom prefixes with \cmd{citemspprefix}\texttt{\{letter\}\{label\}}. Unrecognised first characters fall back to the slot default (\S\ or \P).

  \item \textbf{Citation rendering.} A custom biblatex cite command, \cmd{citelinksp}, declared via \verb|\DeclareCiteCommand|, prints the hyperlinked citation number and conditionally attaches superscript/subscript locators through the dispatcher. In italic contexts, additional math kerns align the locators with the slanted axis. When \pkg{hyperref} draws bordered link boxes (\texttt{colorlinks=false}), a horizontal kern prevents overlap.

  \item \textbf{Entry parsing.} The user-facing \cmd{citemsp} uses \verb|\forcsvlist| (from \pkg{etoolbox}, loaded by \pkg{biblatex}) to iterate over comma-separated entries. Each entry is split at forward slashes using a delimited \verb|\def| pattern, extracting the key and up to two locator slots before calling the renderer.
\end{enumerate}


% ══════════════════════════════════════════════════════════════════════════════
\section{Conventions}
% ══════════════════════════════════════════════════════════════════════════════

\begin{itemize}[leftmargin=2em]
  \item \textbf{Section numbering:} use the source's native format (e.g., 3.2.1, IV, A.1).
  \item \textbf{Paragraph counting:} count paragraphs from the start of the referenced section.
  \item \textbf{Prefix locators:} use a single-letter prefix for typed locators: \texttt{b}~(box), \texttt{d}~(definition), \texttt{e}~(equation), \texttt{f}~(figure), \texttt{n}~(footnote), \texttt{p}~(page), \texttt{t}~(table). E.g., \texttt{key/3.2/e4.3} for an equation, \texttt{key/b21.1} for a box.
  \item \textbf{Subsection references:} use the full path (e.g., \texttt{key/3.2.1}), not a relative number.
  \item \textbf{Plain citations:} \verb|\citemsp{key}| with no slashes produces a standard numeric citation, so the command can replace \verb|\cite| if desired.
\end{itemize}


% ══════════════════════════════════════════════════════════════════════════════
\section{Requirements}
% ══════════════════════════════════════════════════════════════════════════════

\begin{itemize}[leftmargin=2em]
  \item \LaTeX2e\ (2020/10/01 or later)
  \item \pkg{biblatex} (any numeric style) with Biber backend
  \item \pkg{graphicx}
  \item \pkg{amssymb} (for $\square$, $\triangleq$, $\boxdot$, $\boxplus$, $\dagger$ locator glyphs)
\end{itemize}

% ══════════════════════════════════════════════════════════════════════════════
\section{License}
% ══════════════════════════════════════════════════════════════════════════════

This work is released under the \href{https://www.latex-project.org/lppl/lppl-1-3c/}{LaTeX Project Public License, version 1.3c} or later.


\printbibliography

\end{document}
